\cotefile{157-1}
\scanpage{227}\mspage{82}
\[
\Phi\longleftrightarrow\Psi\quad\Longleftrightarrow\quad
\varphi^{-1}(\Phi\times\Psi)\subset\unsure{\mathcal{D}|\mathscr{D}}_0
\]
Mais à vrai dire, je m'aperçois que c'est
quand $=$ \unsure{du dual} d'une relation \deleted{\unsure{entre} $\operatorname{Fl}($}
\unsure{il|qu'il} \deleted{\unsure{s'agit}} \unsure{en|sur} $\operatorname{Fl}(M)$, \deleted{\unsure{puisque}} et cette relation
\deleted{\unsure{peut s'é}} entre $i:A\unsure{-|\to}B$ et $p:X\to Y$
peut s'écrire \unsure{ainsi|comme} :
\[
\begin{aligned}
&\text{pour ttes flèches } A-X,\ B-Y\text{ rendant}\\
&\text{commutatif le diagramme carré, la}\\
&\text{flèche }\varphi\text{ pointillée du carré existe.}
\end{aligned}\tag*{$(**)$}
\]

Cela me fait penser que si on a
une relation $\mathcal{R}$ entre $\unsure{\mathfrak{P}|\mathcal{P}}(P)$ et $\mathfrak{P}(Q)$,
satisfaisant $\unsure{D_2|D|D_1}$, donc décrite par une
application décroissante
\[
\mathfrak{P}(P)\longrightarrow\mathfrak{P}(Q)
\]
\unsure{transformant} $\operatorname{Sup}$ en $\operatorname{Inf}$, \deleted{\unsure{telle}} \inserted{la donnée d'}une
telle application équivaut à celle d'une
application quelconque
\[
P\unsure{-|\to}\mathfrak{P}(Q),
\]
ou ce qui revient au $=$, d'une relation
$R$ sur $P\times Q$. Donc \unsure{cette} structure
\deleted{\unsure{devait}} \unsure{doit|peut} être regardée comme ``de première
espèce''. La situation change à partir
du moment où, pour décrire la relation
$\mathcal{R}$,\inserted{ ou les propriétés voulues,} on se trouve \unsure{contraint} de remplacer
$\mathfrak{P}(P)$, $\mathfrak{P}(Q)$ par des sous-ensembles
convenables -- quand on ne regarde que
\uline{\unsure{certaines}} parties de $P,Q$. C'est ce qu'on
% archive sheet number: 227
